\cotefile{13}
\scanpage{2}
\begin{flushright}
\uline{\unsure{Groupe} de Hodge}\\
\uline{Lettre à Mumford}
\end{flushright}
% archive sheet number: 1

\scanpage{4}
$G$ \unsure{groupe} algébrique sur $\mathbb{Q}$

\[
\unsure{\mathbb{S}|\mathrm{S}}=\prod_{\mathbb{C}/\mathbb{R}}G_{m\,\mathbb{C}},\quad
\text{donc}\quad \mathbb{S}(\mathbb{R})\simeq\mathbb{C}^{*}.\quad
\text{Soit }\underset{\substack{\Vert\\
\operatorname{Ker}(\mathrm{Norme}:\mathbb{S}\to G_{m\,\mathbb{R}})}}{\mathbb{U}}\subset\mathbb{S}
\]

$\operatorname{Rep}(\mathbb{S})\approxeq$ catégorie des vectoriels réels $V_{\mathbb{R}}$, avec bigraduation sur $V_{\mathbb{C}}=V_{\mathbb{R}}\otimes_{\mathbb{R}}\mathbb{C}$,
\[
V_{\mathbb{C}}\simeq\coprod_{p,q}V_{\deleted{\unsure{p,q|\mathbb{C}}}}^{p,q}
\]
\margin{\uline{NB}\quad $\lambda\in\mathbb{C}^{*}\simeq\mathbb{S}(\mathbb{R})$ opère sur $V^{pq}$ par $\lambda^{p}\bar{\lambda}^{q}$}
satisfaisant la condition
\[
V^{p,q}=\overline{V^{q,p}}.
\]

L'hom canonique $G_{m\,\mathbb{R}}\xrightarrow{i}\mathbb{S}$ correspond à la graduation totale

\diagram{\begin{tikzcd}[ampersand replacement=\&]
\operatorname{Rep}(G) \arrow[r,"\text{oubli}"] \&
\operatorname{Modf}(\mathbb{Q}) \arrow[r,"\otimes_{\mathbb{Q}}\mathbb{R}"] \&
\operatorname{Modf}(\mathbb{R}) \arrow[r,"\otimes_{\mathbb{R}}\mathbb{C}"] \&
\operatorname{Modf}(\mathbb{C})
\end{tikzcd}}

La donnée sur le foncteur composé d'une bigraduation
\[
V_{\mathbb{C}}\simeq\coprod_{p,q}V_{\deleted{\unsure{\mathbb{C}|C}}}^{p,q}
\]
\margin{i.e. d'une \unsure{préstructure} de Hodge, compatible avec la structure sur $\mathbb{Q}$ donnée, et fonctorielle en l'objet de $\operatorname{Rep}(G)$}
satisfaisant
\[
V^{p,q}=\overline{V^{q,p}}
\]
équivaut à la donnée d'un hom
\[
\mathbb{S}\xrightarrow{\unsure{j|s}}G_{\mathbb{R}}
\]
Si on se donne un hom [central]
\[
i:G_m\unsure{\rightharpoonup|\to}G\unsure{.}\;,
\]
i.e. une graduation :
\[
V=\coprod_n V^n
\]
sur le foncteur oubli \unsure{$w$|$\omega$}, alors la condition
\[
V^n_{\mathbb{C}}=\coprod_{p+q=n}V^{p,q}
\]
% archive sheet number: 3

\scanpage{5}
équivaut \unsure{à|a} la condition
\[
j\mid G_{m\,\mathbb{R}}=i_{\mathbb{R}}\tag*{(a)}
\]

Pour la suite, $i$ est supposé donné, \deleted{\unsure{Considérons} \unsure{un} $j$ compatible avec $i$ au sens précédent.}\inserted{de $\overline{m}$ qu'un} objet inversible \inserted{$\mathbb{Q}(-1)$} de $\operatorname{Rep}(G)$ à \unsure{isom|iso} près, i.e. un hom \unsure{$\varepsilon$}
\[
\varepsilon:G\longrightarrow G_m
\]
On a alors, pour un $j$ donné, que $\mathbb{Q}(1)$ est de bidegré $(1,1)$ si et seulement si on a les conditions
\[
\varepsilon j(\lambda)=\lambda^2\quad\text{si }\lambda\in\mathbb{R}^{*},\quad
\deleted{\unsure{j(u)=1|j\mathbb{U}=1}}\quad
\varepsilon(j(u))=1\quad\text{si }u\in\mathbb{U}\tag*{(*)}
\]
i.e.
\[
\varepsilon j\deleted{\unsure{\mathbb{S}|j}}=\mathrm{Norme}:\unsure{\mathbb{S}|\mathrm{S}}\longrightarrow G_{m\,\mathbb{R}}\tag*{(b)}
\]
On suppose fixé $i$ et $\varepsilon$ satisfaisant \deleted{\unsure{$(a)$}}\inserted{$\varepsilon i(\lambda)=\lambda^2$}, et on \unsure{s'intéresse} aux $j$ satisfaisant a) et b).

\inserted{Soit $j$ donné.} Cherchons, pour $V\in\operatorname{Ob}\operatorname{Rep}(G)$ donné\inserted{ de degré $p$}, \unsure{d'où|D'où} structure de Hodge sur $V_{\mathbb{C}}$, (ou sur $V_{\mathbb{R}}$) quelles sont les \deleted{\unsure{polarisations} de cette structure de Hodge, \unsure{provenant}}\inserted{formes de $\unsure{R(G)|R(\mathcal{G})}$}
\[
V\times V\xrightarrow{\varphi}\mathbb{Q}(-p)
\]
qui définissent une \uline{polarisation} des structures de Hodge correspondantes. Ce sont celles qui satisfont \unsure{sur|une} la condition que
\[
\varphi_{\mathbb{R}}(x,Cy)\quad\text{forme \inserted{\uline{symétrique}} \uline{définie positive sur} } V_{\mathbb{R}}\times V_{\mathbb{R}}\tag*{(**)}
\]
% archive sheet number: 4

\scanpage{6}
où $C\in\unsure{\mathbb{S}|\mathrm{S}}(\mathbb{R})\simeq\mathbb{C}^{*}$ correspond à l'\unsure{élément} $\sqrt{-1}$ de $\mathbb{C}^{*}$.

[Ceci implique que $\varphi$ est alternée si $n$ impair, symétrique si $n$ pair]
\margin{\uline{NB} Soit $G'=\operatorname{Ker}\varepsilon$, \deleted{\unsure{On a} $\unsure{j(G_m)}$}, alors $G$ est \deleted{\unsure{un quotient de}} $\simeq G'\mathop{\times}\limits_{\mathbb{Z}/2\mathbb{Z}}G_{m\,\mathbb{Q}}$, où $\mathbb{Z}/2\mathbb{Z}\longrightarrow G'$ est induit par $\unsure{i|\varepsilon}$. Donc $G$ (avec $i,\varepsilon$) se \unsure{récupère} en connaissant $G'$ et son point d'ordre 2 $i(-1)=\unsure{\eta|z}$. La donnée de \tikz[baseline=-.25ex,x=1ex,y=1ex]{\draw[line width=.5pt] (0,2.1) -- (0,.18) .. controls (0,-.12) and (.63,-.18) .. (.48,.12) .. controls (.40,.28) and (.18,.23) .. (0,.14);}}

\uline{Théorème} Pour $j$ donné, \unsure{relié|relé} à $i$ et $\varepsilon$ par condition a) et b), les conditions suivantes \inserted{(i)(ii)(i bis)(ii bis)} sont équivalentes et impliquent (iii) :

(i) Pour tout $V\in\operatorname{Ob}\operatorname{Rep}(G)$, \inserted{de degré $n$,} il existe une forme $\varphi:V\deleted{\unsure{\otimes}}\otimes V\longrightarrow\mathbb{Q}(-n)$ dans $\operatorname{Rep}(G)$ qui \unsure{définit} une polarisation (pour \deleted{\unsure{les}}la structure de Hodge \unsure{associée|associées} à $j$ sur $V$, \unsure{$j$|$]$|$\}$} i.e. en fait : $C=j(i)$]

(i bis) Comme (i), mais \deleted{\unsure{avec}} pour des $V_{\alpha}$ donnés seulement \inserted{comme \unsure{module|modules} sur $G'$}, \deleted{qui \unsure{engendrent}}\inserted{et} tels que \deleted{\unsure{les}} l'hom $G\longrightarrow\underline{\mathrm{Aut}}_{\mathbb{Q}}(V_{\alpha})$ soit injectif. De plus, on suppose $G'$ réductif i.e. $G$ réductif

(ii), \deleted{\unsure{(ii bis) Le centralisateur de}} Comme (i), (i bis), mais avec $G$ remplacé par $G_{\mathbb{R}}$, et $\operatorname{Rep}(G)$ par $\operatorname{Rep}(G_{\mathbb{R}})$.

(iii) \inserted{\deleted{\unsure{$G$}} $G'_{\mathbb{R}}$ est réductif, et si $K\deleted{\unsure{(j)}}$ est} Le centralisateur de $j(C)$ \inserted{(noté encore $C$)} dans $G_{\mathbb{R}}^{\prime\,0}\deleted{\unsure{/i_{\mathbb{R}}(G_{m/\mathbb{R}}^{0})}}$\margin{\unsure{composante|compos. de} \deleted{\unsure{connexe}} topologique}, \inserted{l'image de $K(j)$ dans $G_{\mathbb{R}}^{0}/i_{\mathbb{R}}(\mathbb{R}^{*+})$} est un sous-groupe compact maximal, \deleted{\unsure{le}} i.e. le centralisateur de $C$ dans $G_{\mathbb{R}}^{\prime\,0}$ est un ss-groupe compact \deleted{\unsure{en groupe}}\inserted{maximal.}

% margin note written sideways
\margin{\uline{NB} (iii) équivaut aux conditions :

(iii bis) $G^{\mathrm{ad}}=G/\mathrm{Centre}$ avec le composé \deleted{\unsure{la}} $\mathbb{S}\unsure{\mathbin{-}|\to}G\to G^{\inserted{\mathrm{ad}}}\deleted{\unsure{/\mathrm{Centre}|/G_{\mathrm{cen}}}}$ satisfait à (i)

(iii ter) $\exists$ un quotient \inserted{$G''$} de $G$ par un ss-gpe fini de $G'$, tel que le composé $\mathbb{S}\mathbin{-}G\unsure{\rightharpoonup|\to}G''$ satisfasse à (i)}

\noindent\rule{\linewidth}{0.3pt}

\tikz[baseline=-.25ex,x=1ex,y=1ex]{\draw[line width=.5pt] (.34,2.5) -- (.28,.16) .. controls (-.02,.39) and (-.32,.02) .. (-.10,-.14) .. controls (.11,-.28) and (.29,-.03) .. (.28,.16);}\quad $j:\mathbb{S}\longrightarrow G_{\mathbb{R}}$ \deleted{\unsure{équivaut}} satisfaisant a) et b équivaut à la donnée d'un hom $j':\mathbb{U}\longrightarrow G_{\mathbb{R}}'$ satisfaisant $j'(-1)=\eta$
% archive sheet number: 5

\scanpage{7}
\deleted{Pour tout espace analytique $S$, soit $F(S)$ l'ensemble \inserted{\unsure{quotient} des objets suivants} des \uline{foncteurs} qui, à chaque $V\in\operatorname{Rep}(G)$, associe}

Soit $S$ un espace analytique. Une $G$-\uline{subordination} \inserted{\uline{rigidifiée}} \uline{sur} $S$ est la donnée d'un $\otimes$-foncteur qui, à tout $V\in\operatorname{Ob}\operatorname{Rep}(G)$, associe une structure de famille de Hodge \inserted{mod isogénie} sur $V_{\mathbb{C}}\times S\unsure{\underset{\mathrm{D}}{=}|\underset{\delta}{=}|\underset{\mathrm{Df}}{=}}V_S$ \deleted{\unsure{polarisables|polarisable} \unsure{en} chaque point,} [i.e. une bigraduation \inserted{analytique \uline{réelle}} sur $V_{\mathbb{C}}\times S$, \inserted{\deleted{\unsure{autres}} avec $V_S^{pq}=\overline{V_S^{q,p}}$,} telle que la $p$-filtration correspondante soit une filtration analytique complexe] et telle que \inserted{a) degré total soit celui défini par $i$ b) $\mathbb{Q}(-1)$ devienne \unsure{St|Sf} \unsure{de H.} de bidegré $(1,1)$, et} \inserted{c)} pour \inserted{$\forall$ $V$ homogène de degré $n$ et} $\forall\,s\in S$, la fibre en $s$ soit polarisable \deleted{\unsure{pour} $\forall$ \unsure{point fixé}}\inserted{à l'aide d'une forme $\varphi:V\otimes V\to\mathbb{Q}(-n)$ \inserted{qui soit $G$-invariante (NB. si $S$ est connexe, il suffit de l'exiger cette condition en \uline{un} point \deleted{\unsure{$s$}} de $S$, car la forme restera de polarisation partout, puisqu'elle est \unsure{non dég.} partout}}

\deleted{\unsure{Retenant|Retenons}}\inserted{Oubliant pour le moment la condition de} holomorphie de la filtration, on trouve :

a) pour \unsure{tt} \unsure{pt} $s\in S$, on veut une façon \deleted{de faire} fonctorielle de mettre une structure de Hodge \inserted{(mod isogénie)} polarisable sur les $V\in\operatorname{Ob}\operatorname{Rep}(G)$, i.e. donnée par un $j=\unsure{J|j}(s):\unsure{\mathbb{S}|\mathrm{S}}\longrightarrow G_{\mathbb{R}}$ satisfaisant les conditions du th, i.e. $j'=J'(s):\mathbb{U}\unsure{\rightharpoonup|\to}G'_{\mathbb{R}}$ satisfaisant $j'(\deleted{\unsure{C}}-1)=\unsure{\eta|z}$, et tel que $\unsure{\operatorname{Zent}|\operatorname{Zentr}}_{G_{\mathbb{R}}^{\prime\,0}}(j'(s)(C))$ soit compact maximal.

b) Pour $s$ variable, $j'(s)$ dépend de $s$ de façon analytique réelle.
% archive sheet number: 6

\scanpage{8}
Enfin, pour exprimer la condition d'holomorphie de la filtration, il convient d'introduire, \deleted{\unsure{sur}} sur l'\unsure{ens} \inserted{$\unsure{\mathscr{J}|\mathcal{J}|\gamma}$} des $j$ (ou $j'$) possibles, une structure de \uline{variété} \uline{analytique} \uline{complexe} naturelle, et la condition s'exprime par l'\uline{holomorphie} de l'application
\[
s\longmapsto\unsure{J|j}(s):\quad S\unsure{\rightharpoonup|\longrightarrow}\mathscr{J}.
\]

\deleted{En un \unsure{mot}} Donc l'ens des $G$-subordinations \inserted{rigidifiées} \unsure{base|sur} $S$ s'identifie à l'ens. $\operatorname{Mor}(S,\mathscr{J})$\margin{\unsure{morph.} d'espaces analytiques}. Donc pour $S$ variable, on a un foncteur représentable par $\mathscr{J}$.

\uline{NB} $G(\mathbb{R})$ opère de façon évidente sur $\mathscr{J}$ (mod \unsure{autom|automs} intérieurs) ; les orbites sont \uline{ouvertes}. Donc

\deleted{Soit \deleted{\unsure{$G$ un}} $V$ un $G$-module fidèle, et $V_0\subset V$ un $Z$-module de type fini engendrant $V$. Alors \unsure{tte|tt} $G$-subordination sur un $S$ définit une structure de Hodge (\uline{pas} mod isog) sur $V_0\times S$, Soit $\Gamma\subset G(\mathbb{Q})$ le stabilisateur de $V_0$.}

les comp. connexes de $\mathscr{J}$ sont de la forme $G(\mathbb{R})^{0}/H(\mathbb{R})\cap G(\mathbb{R})^{0}$, où $H$ est le centralisateur d'un $j:\unsure{\mathbb{S}|\mathrm{S}}\longrightarrow G$ dans $G$ ; ou encore $G'(\mathbb{R})^{0}/H'(\mathbb{R})\cap G'(\mathbb{R})^{0}$, $H'$ le centralisateur d'un $j':\mathbb{U}\longrightarrow G'$ dans $G'(\mathbb{R})^{0}$.
% archive sheet number: 7

\scanpage{9}
Soit $C=\operatorname{Rep}G$ la $\otimes$-catégorie définie par $G$, et $S$ un espace analytique \uline{connexe}, $s\in S$. On veut classifier les $\otimes$-foncteurs
\[
F:C\longrightarrow\text{str. de Hodge \inserted{(mod. isogénie)} sur }S
\]
\unsure{telles|tels} que\tikz[baseline=-.2ex,x=1ex,y=1ex]{\draw[line width=.55pt] (0,2.0)--(0,.13) .. controls (-.30,-.07) and (.32,-.16) .. (.17,.10) .. controls (.13,.21) and (.05,.19) .. (0,.13);}\inserted{\tikz[baseline=-.2ex,x=1ex,y=1ex]{\draw[line width=.42pt] (.10,2)--(0,.16) .. controls (-.26,-.02) and (.37,-.20) .. (.28,.11) .. controls (.23,.25) and (.08,.24) .. (0,.16);}\ a) degré total (\unsure{rel.} à $i$) $=\sum$ degrés partiels b) $\mathbb{Q}(-1)$ \unsure{donne|devienne} str. de Hodge de bidegré $(1,1)$, et} c) Pour \deleted{tout $s\in S$ [en \uline{un} $s\in S$, c'est kif kif]} \deleted{et} $V\in\operatorname{Ob}C$, \deleted{\unsure{polarisables en chaque point}} \inserted{$V$ de degré $n$, $\exists\,\varphi:V\otimes V\unsure{\rightharpoonup|\to}\mathbb{Q}(-n)$ tel que $F(\varphi)$ soit une polarisation [en un pt, en \unsure{tt|tout} point, de $S$, c'est kif-kif].}
\margin{De tels objets s'appellent \uline{$G$-subordination sur $S$} (pas \unsure{néc.|néc} \unsure{rigidifiée|rigidifiées}!)}
Composant avec le foncteur \unsure{fibre} en $s$, on trouve un foncteur fibre
\[
C\longrightarrow\operatorname{Modf}(\mathbb{Q}),
\]
et quitte à tordre $G$, on peut supposer qu'on se donné un isomorphisme \inserted{$\varphi$} entre ce dernier et le foncteur oubli.

\deleted{Alors} Soit $\pi=\pi_1(S,\unsure{\xi|\zeta|s})$. Alors $\pi$ opère \inserted{à g.}, sur le foncteur \deleted{\unsure{``fibre}} ``$\mathbb{Q}$-fibre en $s$'' sur la catégorie des str. de Hodge sur $S$, donc opère sur le foncteur composé avec $F$, d'où un premier invariant de $(F,\varphi)$ :
\[
\alpha:\pi=\pi_1(S,\xi)\longrightarrow G(\mathbb{Q})
\]
Soit d'autre part $\widetilde{S}$ le \unsure{revêtt|revêtmt} universel de $S$ \deleted{\unsure{pointé}} \unsure{marqué} sur $s$ par $\widetilde{s}$. Alors la donnée de $\varphi$ permet de considérer le foncteur $V\longmapsto F(V)\longmapsto F(\unsure{W|V})\times_{\unsure{s|S}}\overline{S}$ comme une $G$-subordination de $\widetilde{S}$. Donc on trouve un
% archive sheet number: 8

\scanpage{10}
\unsure{deuxième} élément de structure
\[
f:\widetilde{S}\longrightarrow\unsure{\mathscr{J}|\mathcal{J}|\gamma}
\]
Ces deux éléments de structure sont reliés par la condition que $f$ est compatible avec les actions de $\pi$, quand on le fait agir sur $\mathscr{J}$ via $\alpha:\pi\unsure{\dashv|\to}G(\mathbb{Q})\subset G(\mathbb{R})$ et l'action \inserted{à \uline{droite}} de $G(\mathbb{R})$ sur $\mathscr{J}$ (opposée à l'action gauche).

La donnée d'un couple $(F,\varphi)$ équivaut à celle d'un couple $(f,\alpha)$ satisfaisant les deux conditions précédentes.

\deleted{Il faudrait fixer des \uline{réseaux}, et des sous-groupes discrets convenables de $G(\mathbb{Q})$, du type $G(\mathbb{Z})$ commensurable à un tel \unsure{....|...}}

Autre façon de présenter les \unsure{``}choses, \unsure{indépend.ment|indépendamment} du choix d'un point-base $s$ de $S$. \deleted{(\unsure{reprenons|recommençons})}.

\deleted{On dit que $F$ est \unsure{co|con}} Notons que les foncteurs fibres \inserted{$\unsure{\Phi|\varphi}_s$ (rel/$\mathbb{Q}$)} sur \deleted{str. de Hodge} $\mathrm{Hodge}(S)$ définis par les divers points $s\in S$ \deleted{\unsure{sont}} \inserted{d'une \unsure{$\overline{m}$|\deleted{c}} composante connexe $S^{i}$ de $S$ sont} isomorphes ; donc aussi les $\Phi_s\circ F$ \inserted{$(s\in S^{i})$}. On dit que \inserted{la $G$-subordination} $F$ est \uline{compatible avec le foncteur} \uline{oubli} \uline{sur} $C=\operatorname{Rep}(G)$ si les $\Phi_s\circ F$ sont isomorphes audit foncteur oubli. Soit alors, pour
% archive sheet number: 9

